\documentclass[10pt,a4paper]{article}
\usepackage[margin=2.2cm]{geometry}
\usepackage{amsmath,amssymb,booktabs,graphicx,microtype,siunitx}
\usepackage[table]{xcolor}
\usepackage[hidelinks]{hyperref}
\usepackage{caption}
\usepackage{times}
\usepackage[section]{placeins}
\captionsetup{font=small,labelfont=bf}
\sisetup{detect-all}
\input{macros.tex}

\title{\vspace{-1.3cm}\textbf{Routing, capacity and width in dynamic convolution\\ across five backbones}\\[3pt]
\large Cross-model summary of the replicated 132-run program}
\author{Ankit Kumar\\\small Mechatronics Department, University of Agder, Grimstad, Norway}
\date{\small Generated from result files by \texttt{report/cross.py}}

\begin{document}
\maketitle
\vspace{-0.4cm}
\V{synthbanner}

\noindent The 132-run ResNet-18 program (ten groups, 22 arms, six seeds)
was repeated on each backbone with the same recipe, arms, seeds, margin
and dynamic-layer rule. This summary puts \V{nmodels} backbones side by
side (\V{nruns} runs, \V{gpuh} GPU-hours logged). Every number is
computed from the result files by the code that builds the per-model
reports, so the two always agree. Backbones that did not pass the gate or
have no results are listed in Table~\ref{tab:xsummary}, not dropped.

\section{The backbones}

\begin{table}[h]\centering\small
\caption{Per backbone, CIFAR-100 width 64: stored parameters of the static
and $K=4$ \texttt{gap} networks, the share of the static network's
parameters that sit in dynamic layers, the \texttt{gap} router's size
relative to the static network, the gate run, the static mean over
six seeds, runs present, diverged runs, groups that needed gradient
accumulation, and logged GPU-hours.}
\label{tab:xsummary}
\resizebox{\textwidth}{!}{%
\begin{tabular}{l rr r r rr r r r r}
\toprule
backbone & static (M) & gap (M) & dyn.\ share (\%) & router (\% of static) & gate (\%) & static mean (\%) & runs & diverged & accum.\ groups & GPU-h \\
\midrule
%%XSUMMARY%%
\bottomrule
\end{tabular}}
\end{table}

\section{Bank against routing}

\begin{figure}[h]\centering
\includegraphics[width=\textwidth]{fig/x1_decomp.pdf}
\caption{(a) Bank term, \texttt{uniform}$-$\texttt{static}, and (b)
routing term, \texttt{gap}$-$\texttt{uniform}, for every backbone and
condition (C: CIFAR-100, T: Tiny-ImageNet, number: width), \SI{95}{\percent}
confidence intervals.}
\label{fig:xdecomp}
\end{figure}

The bank term is equivalent to zero (TOST, $\pm$\SI{0.5}{pp}) in
\V{bankeq} of \V{bankn} backbone--condition cells, and the routing term
is positive with $p<0.05$ in \V{routpos} of \V{bankn}
(Tables~\ref{tab:xbank},~\ref{tab:xrout}).

\begin{table}[h]\centering\small
\caption{Bank term (pp); $\equiv$ marks equivalence to zero by TOST.}
\label{tab:xbank}
\begin{tabular}{l rrrrr}
\toprule
backbone & C16 & C32 & C64 & T32 & T64 \\
\midrule
%%XBANK%%
\bottomrule
\end{tabular}
\end{table}

\begin{table}[h]\centering\small
\caption{Routing term (pp); * marks $p<0.05$.}
\label{tab:xrout}
\begin{tabular}{l rrrrr}
\toprule
backbone & C16 & C32 & C64 & T32 & T64 \\
\midrule
%%XROUT%%
\bottomrule
\end{tabular}
\end{table}

\section{Total effect by width}

\begin{figure}[h]\centering
\includegraphics[width=\textwidth]{fig/x2_total.pdf}
\caption{\texttt{gap}$-$\texttt{static} against width, every backbone,
\SI{95}{\percent} CI; the band is the margin.}
\label{fig:xtotal}
\end{figure}

\begin{table}[h]\centering\small
\caption{\texttt{gap}$-$\texttt{static} (pp) and verdict per condition.}
\label{tab:xtotal}
\resizebox{\textwidth}{!}{%
\begin{tabular}{l lllll}
\toprule
backbone & C16 & C32 & C64 & T32 & T64 \\
\midrule
%%XTOTAL%%
\bottomrule
\end{tabular}}
\end{table}

\section{Equal parameter budgets}

\begin{figure}[h]\centering
\includegraphics[width=\textwidth]{fig/x3_matched.pdf}
\caption{(a) \texttt{static} at width $2c$ minus \texttt{gap} at width
$c$, \SI{95}{\percent} CI; the number under each bar is the parameter ratio
routed/plain, red where it is outside $\pm$10\,\%. (b) The static accuracy
curve interpolated in log-parameters at the \texttt{gap} network's count,
minus \texttt{gap}; missing points fall outside the static range.}
\label{fig:xmatched}
\end{figure}

Width doubling is a matched-budget comparison only where nearly all
parameters sit in dynamic layers. It is matched (within 10\,\%) in
\V{nmatched} of \V{npairs} backbone--pair cells
(Table~\ref{tab:xmatched}); elsewhere the table gives the ratio and the
interpolated equal-parameter comparison.

\begin{table}[h]\centering\small
\caption{\texttt{static}($2c$) $-$ \texttt{gap}($c$), parameter ratio and
the interpolated comparison at equal parameters (no CI).}
\label{tab:xmatched}
\resizebox{\textwidth}{!}{%
\begin{tabular}{l l r l l r l r}
\toprule
backbone & pair & $\Delta$ (pp) & \SI{95}{\percent} CI & $p$ & ratio & label & $\Delta$ at equal params \\
\midrule
%%XMATCHED%%
\bottomrule
\end{tabular}}
\end{table}

\begin{figure}[h]\centering
\includegraphics[width=\textwidth]{fig/x5_curves.pdf}
\caption{CIFAR-100 accuracy against stored parameters, widths 16, 32 and
64, per backbone.}
\label{fig:xcurves}
\end{figure}

\section{Temperature, bank size and routing entropy}

\begin{figure}[h]\centering
\includegraphics[width=\textwidth]{fig/x6_temp_k.pdf}
\caption{(a) Annealed minus constant temperature, \texttt{gap},
CIFAR-100. (b) $K=8$ minus $K=2$, CIFAR-100 width 32.}
\label{fig:xtempk}
\end{figure}

\begin{table}[h]\centering\small
\caption{Annealed $-$ constant (pp, $p$) per width, and $K=8-K=2$.}
\label{tab:xtempk}
\resizebox{\textwidth}{!}{%
\begin{tabular}{l lll r l}
\toprule
backbone & w16 & w32 & w64 & $K=8-K=2$ & verdict \\
\midrule
%%XTEMPK%%
\bottomrule
\end{tabular}}
\end{table}

\begin{figure}[h]\centering
\includegraphics[width=0.62\textwidth]{fig/x4_entropy.pdf}
\caption{Router flatness $1-\tilde H$ against the benefit over
\texttt{uniform}, CIFAR-100, three widths and two recipes per backbone;
$\rho$ is Spearman's correlation within a backbone.}
\label{fig:xent}
\end{figure}

\begin{table}[h]\centering\small
\caption{Entropy against benefit: points, Spearman $\rho$, and at width 16
the ratio of constant to annealed $1-\tilde H$.}
\label{tab:xent}
\begin{tabular}{l r r r}
\toprule
backbone & points & $\rho$ & flatness ratio (w16) \\
\midrule
%%XENT%%
\bottomrule
\end{tabular}
\end{table}

\section{Pre-registered predictions}

Recorded in \texttt{PREREGISTRATION\_MODELS.md} on 30 September 2026 for
the \V{nnew} new backbones before any of their runs. P1: the bank term is
equivalent to zero in all five conditions (held for \V{p1held} of
\V{nnew}). P2: \texttt{static} at width $2c$ beats \texttt{gap} at width
$c$ in all three pairs (held for \V{p2held}). P3: the routing term is
positive with $p<0.05$ in at least three conditions (held for
\V{p3held}). ResNet-18 is shown for reference; for it these are post hoc.

\begin{table}[h]\centering\small
\caption{Pre-registered predictions per backbone.}
\label{tab:xprereg}
\begin{tabular}{l rl rl rl}
\toprule
& \multicolumn{2}{c}{P1 bank $\equiv 0$} & \multicolumn{2}{c}{P2 width wins} & \multicolumn{2}{c}{P3 routing $>0$} \\
\cmidrule(lr){2-3}\cmidrule(lr){4-5}\cmidrule(lr){6-7}
backbone & count & outcome & count & outcome & count & outcome \\
\midrule
%%XPREREG%%
\bottomrule
\end{tabular}
\end{table}

\section{Threats to validity}

\paragraph{One recipe for five architectures.} Every backbone is trained
with the ResNet-18 recipe (SGD, lr 0.1, 150 epochs). That is the design:
a per-architecture recipe would change two things at once. It also means
absolute accuracies, above all ConvNeXt-Tiny's, sit below what a tuned
recipe reaches. The gate guards against a backbone the recipe does not
train at all; it does not make the recipes equally good.

\paragraph{Placement.} The ResNet-18 rule makes different shares of each
network dynamic (Table~\ref{tab:xsummary}): nearly all of ResNet-18 and
VGG-16, a smaller part of ResNet-50, and only depthwise layers in
MobileNetV2 and ConvNeXt-Tiny. A backbone-to-backbone difference may come
from this share rather than from the architecture.

\paragraph{Router size.} Every backbone uses the ResNet-18 router, a
two-layer MLP of hidden width $C_{\text{in}}/4$. Its size grows with the
square of the input channels while a depthwise kernel grows linearly, so
in the depthwise backbones the router is much larger than the bank it
routes (Table~\ref{tab:xsummary}). The routing term
\texttt{gap}$-$\texttt{uniform} then also compares networks of very
different stored size; \texttt{frozen\_random} (an untrained router,
held as buffers) is the control for this at CIFAR-100 width 64.

\paragraph{Initialisation.} In the 132 ResNet-18 runs static convolutions
kept PyTorch's default initialisation while bank kernels were
Kaiming-normal. The new backbones use one scheme for both. The routing
term (\texttt{gap}$-$\texttt{uniform}) is unaffected in either case,
since both arms have banks; the bank term of ResNet-18 compares two
initialisations as well as two architectures.

\paragraph{Gradient accumulation.} Where batch 128 did not fit, a group
ran with micro-batches and accumulation (Table~\ref{tab:xsummary}). The
optimiser steps and schedule are unchanged; BatchNorm statistics are per
micro-batch. The choice is per group, so every paired comparison within a
group shares it.

\end{document}
